\documentclass[11pt]{article}

\usepackage[margin=1in]{geometry}
\usepackage{amsmath, amssymb, amsthm}
\usepackage{bm}
\usepackage{mathtools}
\usepackage{enumitem}
\usepackage{xcolor}
\usepackage{hyperref}
\usepackage{listings}
\usepackage{booktabs}

\hypersetup{
    colorlinks=true,
    linkcolor=blue,
    urlcolor=blue,
    citecolor=blue
}

\definecolor{codegray}{gray}{0.95}
\lstset{
    backgroundcolor=\color{codegray},
    basicstyle=\ttfamily\small,
    frame=single,
    breaklines=true,
    columns=fullflexible
}

\title{Policy Gradient Reinforcement Learning on CartPole\\
\large A Fundamental Study Note for Implementation}
\author{}
\date{}

\begin{document}

\maketitle

\section{Goal of This Note}

We want to understand how to train a neural network policy on the CartPole environment.
CartPole is a classic reinforcement learning example where a cart moves left or right in order to keep a pole balanced.

The important idea is:

\[
\boxed{
\text{state} \rightarrow \text{neural network policy} \rightarrow \text{action}
}
\]

The neural network receives the current state of the environment and outputs a probability distribution over actions.
Then we sample or choose an action, apply it in the environment, receive reward, and use policy gradient to make good actions more likely in the future.

\section{Prerequisite Concepts}

Before implementing CartPole with policy gradient, we need a few concepts.

\subsection{State}

A state is the information the environment gives to the agent at a time step.
For CartPole, the state is a vector with four real numbers:

\[
 s_t =
 \begin{bmatrix}
 x_t \\
 \dot{x}_t \\
 \theta_t \\
 \dot{\theta}_t
 \end{bmatrix}
\]

where:

\begin{itemize}[leftmargin=2em]
    \item $x_t$ is the cart position;
    \item $\dot{x}_t$ is the cart velocity;
    \item $\theta_t$ is the pole angle;
    \item $\dot{\theta}_t$ is the pole angular velocity.
\end{itemize}

So CartPole gives the agent a four-dimensional observation:

\[
\boxed{s_t \in \mathbb{R}^4}
\]

\subsection{Action}

An action is what the agent chooses to do.
In CartPole, there are two possible actions:

\[
 a_t \in \{0,1\}
\]

Usually:

\begin{itemize}[leftmargin=2em]
    \item $a_t = 0$: push the cart left;
    \item $a_t = 1$: push the cart right.
\end{itemize}

So CartPole has a discrete action space with two actions:

\[
\boxed{|\mathcal{A}| = 2}
\]

\subsection{Reward}

At each step, the environment returns a reward.
In CartPole, the agent usually receives reward $+1$ for every time step that the pole remains balanced.

So if the pole stays balanced for $200$ steps, the total reward is approximately:

\[
 R_{\text{episode}} = 200
\]

The agent's goal is to maximize total future reward.

\subsection{Episode}

An episode is one full run of the environment.
For CartPole, an episode starts with the pole almost upright and ends when the pole falls too much, the cart moves too far, or the maximum time limit is reached.

A trajectory, or rollout, is the sequence:

\[
\tau = (s_0, a_0, r_0, s_1, a_1, r_1, \dots, s_T)
\]

This is the data generated by one episode.

\section{The Policy}

The policy is the agent's behavior rule.
It tells us which action the agent should take in each state.

In stochastic reinforcement learning, the policy is written as:

\[
\pi_\theta(a_t \mid s_t)
\]

This means:

\[
\boxed{
\text{probability of taking action } a_t \text{ given state } s_t
}
\]

The parameter $\theta$ represents the weights of the neural network.

\subsection{Neural Network as Policy}

For CartPole, the neural network receives four inputs and outputs two logits:

\[
\mathbb{R}^4 \rightarrow \mathbb{R}^2
\]

More explicitly:

\[
\begin{bmatrix}
 x_t \\
 \dot{x}_t \\
 \theta_t \\
 \dot{\theta}_t
\end{bmatrix}
\xrightarrow{\text{MLP}_\theta}
\begin{bmatrix}
 z_{\text{left}} \\
 z_{\text{right}}
\end{bmatrix}
\]

The values $z_{\text{left}}$ and $z_{\text{right}}$ are called logits.
They are not probabilities yet.
They are raw scores.

To turn logits into probabilities, we use softmax:

\[
\pi_\theta(a=i \mid s_t)
= \frac{e^{z_i}}{\sum_{j=1}^{2} e^{z_j}}
\]

So the network gives something like:

\[
\pi_\theta(\text{left} \mid s_t) = 0.35,
\qquad
\pi_\theta(\text{right} \mid s_t) = 0.65
\]

Then we can sample an action from this distribution.

\section{Interaction Loop}

At each time step, the agent does the following:

\begin{enumerate}[leftmargin=2em]
    \item Receive current state $s_t$ from the environment.
    \item Pass $s_t$ through the policy neural network.
    \item Get action probabilities $\pi_\theta(a_t \mid s_t)$.
    \item Sample an action $a_t$.
    \item Apply the action in the environment.
    \item Receive reward $r_t$ and next state $s_{t+1}$.
    \item Store the log-probability and reward for training.
\end{enumerate}

Mathematically:

\[
 s_t \rightarrow \pi_\theta(a_t \mid s_t) \rightarrow a_t \rightarrow r_t, s_{t+1}
\]

Implementation-style pseudocode:

\begin{lstlisting}[language=Python]
state = env.reset()
done = False

log_probs = []
rewards = []

while not done:
    logits = policy_net(state)
    probs = softmax(logits)
    action = sample_from(probs)

    log_prob = log_probability_of_chosen_action(probs, action)

    next_state, reward, done, info = env.step(action)

    log_probs.append(log_prob)
    rewards.append(reward)

    state = next_state
\end{lstlisting}

The key thing we store is not just the action.
We store:

\[
\log \pi_\theta(a_t \mid s_t)
\]

This is needed because policy gradient changes the neural network weights using the log-probability of the sampled action.

\section{Return}

The return is the total future reward after a time step.
For time step $t$, the discounted return is:

\[
G_t = r_t + \gamma r_{t+1} + \gamma^2 r_{t+2} + \gamma^3 r_{t+3} + \cdots
\]

Compactly:

\[
\boxed{
G_t = \sum_{k=0}^{T-t} \gamma^k r_{t+k}
}
\]

where $\gamma$ is the discount factor.

\subsection{Discount Factor}

The discount factor $\gamma$ controls how much the agent cares about future rewards.

\[
0 \leq \gamma \leq 1
\]

If $\gamma = 0$, the agent only cares about immediate reward.
If $\gamma$ is close to $1$, the agent cares a lot about future reward.

For CartPole, common choices are:

\[
\gamma = 0.99
\]

This means future rewards are very important, which makes sense because the goal is to keep the pole balanced for as long as possible.

\section{Objective Function}

The policy's goal is to maximize expected return:

\[
J(\theta) = \mathbb{E}_{\tau \sim \pi_\theta}[G(\tau)]
\]

This means:

\[
\boxed{
\text{Choose neural network weights } \theta \text{ that produce high-reward trajectories.}
}
\]

The difficulty is that the action is sampled.
We cannot directly backpropagate through the environment, because the environment may be non-differentiable.
Instead, policy gradient gives us a way to update the policy using sampled experience.

\section{Policy Gradient Core Idea}

The policy gradient tells us how to change $\theta$ to increase expected reward.

The main formula is:

\[
\boxed{
\nabla_\theta J(\theta)
= \mathbb{E}\left[
\nabla_\theta \log \pi_\theta(a_t \mid s_t) G_t
\right]
}
\]

The intuition is very important:

\begin{itemize}[leftmargin=2em]
    \item If an action was followed by high return, increase its probability.
    \item If an action was followed by low return, do not reinforce it much.
    \item If using an advantage estimate, bad actions can be actively pushed down.
\end{itemize}

So if the robot/cart chooses an action that helps keep the pole balanced, the probability of that action in similar states should increase.

\section{Policy Gradient Loss}

PyTorch minimizes losses.
But reinforcement learning wants to maximize reward.
So we use a negative sign.

The REINFORCE-style loss is:

\[
\boxed{
\mathcal{L}(\theta)
= - \sum_{t=0}^{T} \log \pi_\theta(a_t \mid s_t) G_t
}
\]

This is the central implementation formula.

Each chosen action contributes:

\[
-\log \pi_\theta(a_t \mid s_t) G_t
\]

If $G_t$ is large, minimizing this loss increases the probability of action $a_t$ in state $s_t$.

\subsection{Why the Negative Sign?}

If the return is good, we want to maximize:

\[
\log \pi_\theta(a_t \mid s_t) G_t
\]

But optimizers like Adam usually minimize.
So we minimize the negative:

\[
-\log \pi_\theta(a_t \mid s_t) G_t
\]

Thus:

\[
\text{maximize reward-weighted log-probability}
\quad \Longleftrightarrow \quad
\text{minimize negative reward-weighted log-probability}
\]

\section{Simple CartPole Training Algorithm}

The basic algorithm is:

\begin{enumerate}[leftmargin=2em]
    \item Initialize policy network $\pi_\theta$.
    \item Run one episode using the current policy.
    \item Store rewards and log-probabilities.
    \item Compute discounted returns $G_t$.
    \item Compute loss:
    \[
    \mathcal{L}(\theta)
    = - \sum_t \log \pi_\theta(a_t \mid s_t)G_t
    \]
    \item Backpropagate the loss.
    \item Update $\theta$ using an optimizer.
    \item Repeat for many episodes.
\end{enumerate}

Implementation-style pseudocode:

\begin{lstlisting}[language=Python]
for episode in range(num_episodes):
    state = env.reset()
    done = False

    log_probs = []
    rewards = []

    while not done:
        logits = policy_net(state)
        dist = Categorical(logits=logits)
        action = dist.sample()

        log_prob = dist.log_prob(action)

        next_state, reward, done, info = env.step(action.item())

        log_probs.append(log_prob)
        rewards.append(reward)

        state = next_state

    returns = compute_discounted_returns(rewards, gamma=0.99)

    loss = 0
    for log_prob, G in zip(log_probs, returns):
        loss += -log_prob * G

    optimizer.zero_grad()
    loss.backward()
    optimizer.step()
\end{lstlisting}

\section{Why We Sample Instead of Taking the Best Action}

During training, we usually sample actions from the policy distribution:

\[
 a_t \sim \pi_\theta(\cdot \mid s_t)
\]

This allows exploration.

If the agent always takes the most likely action, it may get stuck doing a bad behavior too early.
Sampling lets the policy try different actions and discover better ones.

After training, we may use a deterministic action:

\[
 a_t = \arg\max_a \pi_\theta(a \mid s_t)
\]

That means we choose the most likely action.

\section{On-Policy Nature of This Algorithm}

This basic policy gradient method is on-policy.
That means the data used for training comes from the same policy that is being updated.

At episode $k$, we collect data using the current policy:

\[
\pi_{\theta_k}
\]

Then we update the same policy using that data:

\[
\theta_k \rightarrow \theta_{k+1}
\]

After updating, the old data is no longer exactly from the current policy.
So simple policy gradient methods usually collect fresh data again.

\section{Actor and Critic}

The simple REINFORCE algorithm only uses a policy network.
But more advanced methods often use two networks:

\begin{center}
\begin{tabular}{lll}
\toprule
Network & Formula & Purpose \\
\midrule
Actor / Policy & $\pi_\theta(a \mid s)$ & Chooses actions \\
Critic / Value & $V_\phi(s)$ & Estimates how good a state is \\
\bottomrule
\end{tabular}
\end{center}

The policy is the network we actually use to act.
The critic helps training by reducing variance.

\subsection{Value Function}

The value function estimates expected future reward from a state:

\[
V_\phi(s_t) \approx \mathbb{E}[G_t \mid s_t]
\]

In words:

\[
\boxed{
V_\phi(s_t) = \text{how good this state is expected to be}
}
\]

\subsection{Advantage}

Instead of using raw return $G_t$, we often use advantage:

\[
A_t = G_t - V_\phi(s_t)
\]

Advantage means:

\[
\boxed{
\text{Was this action better or worse than expected?}
}
\]

If:

\[
A_t > 0
\]

then the action was better than expected, so we increase its probability.

If:

\[
A_t < 0
\]

then the action was worse than expected, so we decrease its probability.

The policy loss becomes:

\[
\boxed{
\mathcal{L}_{\text{policy}}(\theta)
= -\sum_t \log \pi_\theta(a_t \mid s_t) A_t
}
\]

This is closer to what algorithms like A2C and PPO use.

\section{CartPole Neural Network Architecture}

A simple policy network for CartPole can be an MLP:

\[
\mathbb{R}^4 \rightarrow \mathbb{R}^{128} \rightarrow \mathbb{R}^{128} \rightarrow \mathbb{R}^{2}
\]

For example:

\[
\begin{aligned}
 h_1 &= \text{ReLU}(W_1s + b_1) \\
 h_2 &= \text{ReLU}(W_2h_1 + b_2) \\
 z &= W_3h_2 + b_3
\end{aligned}
\]

where:

\[
 z \in \mathbb{R}^2
\]

These two output values are logits for left and right.

Then:

\[
\pi_\theta(a \mid s) = \text{softmax}(z)
\]

\section{Implementation Notes}

\subsection{Use Logits Directly}

In PyTorch, it is often better to pass logits directly into a categorical distribution:

\begin{lstlisting}[language=Python]
dist = torch.distributions.Categorical(logits=logits)
action = dist.sample()
log_prob = dist.log_prob(action)
\end{lstlisting}

This is numerically cleaner than manually doing softmax and log.

\subsection{Normalize Returns}

For stability, it is common to normalize returns inside each batch:

\[
\hat{G}_t = \frac{G_t - \mu_G}{\sigma_G + \epsilon}
\]

This does not change the basic idea, but it can make learning more stable.

\begin{lstlisting}[language=Python]
returns = (returns - returns.mean()) / (returns.std() + 1e-8)
\end{lstlisting}

\subsection{Batching Episodes}

One episode can be noisy.
A better implementation can collect multiple episodes before doing one update.

Instead of:

\[
\text{one episode} \rightarrow \text{one update}
\]

we can do:

\[
\text{many episodes} \rightarrow \text{one update}
\]

This gives a better estimate of the policy gradient.

\section{Important Intuition}

Policy gradient is not supervised learning in the usual sense.
There is no label saying:

\[
\text{correct action} = \text{right}
\]

Instead, the environment only says:

\[
\text{this sequence of actions produced this much reward}
\]

Then policy gradient assigns credit to the actions by using:

\[
\log \pi_\theta(a_t \mid s_t)G_t
\]

So the model slowly learns:

\[
\boxed{
\text{In states like this, actions like that tend to produce high future reward.}
}
\]

\section{Full Conceptual Summary}

For CartPole:

\begin{itemize}[leftmargin=2em]
    \item The environment state has four numbers.
    \item The action space has two actions: left or right.
    \item The neural network maps state to two logits.
    \item Softmax turns logits into action probabilities.
    \item The agent samples an action from the policy.
    \item The environment gives reward and next state.
    \item After an episode, we compute discounted returns.
    \item We compute the policy gradient loss.
    \item Backpropagation updates the neural network weights.
    \item Over many episodes, the policy learns to balance the pole.
\end{itemize}

The essential formula is:

\[
\boxed{
\mathcal{L}(\theta)
= -\sum_t \log \pi_\theta(a_t \mid s_t)G_t
}
\]

And the essential intuition is:

\[
\boxed{
\text{Actions that led to high reward become more likely.}
}
\]

\section{Connection to LLM RL}

The same idea applies to language models.
For an LLM, the policy is:

\[
\pi_\theta(\text{next token} \mid \text{previous tokens})
\]

For CartPole:

\[
\pi_\theta(\text{action} \mid \text{state})
\]

The structure is very similar:

\begin{center}
\begin{tabular}{lll}
\toprule
Setting & State / Context & Action \\
\midrule
CartPole & cart position, velocity, pole angle & left or right \\
LLM & previous tokens / prompt & next token \\
\bottomrule
\end{tabular}
\end{center}

In both cases, policy gradient increases the probability of actions that lead to higher reward.

\section{Minimal Mental Model}

The shortest way to remember the whole thing is:

\[
\boxed{
\text{Policy gradient} = \text{reward-weighted cross entropy on sampled actions}
}
\]

But unlike supervised learning, the target action was not given by a dataset.
The action was sampled by the policy itself, and the reward decides whether that sampled action should become more or less likely.

\end{document}
